% AUTO-GENERATED FILE. Source: pub/claims.toml
Evaluating explainability methods requires more than a single faithfulness proxy. We present a modular benchmarking framework centered on quantitative XAI quality metrics: fidelity, stability, sparsity, computational cost, and faithfulness gap, plus an explicit method for operating the framework end-to-end. On the UCI Adult benchmark, we use a staged evidence protocol: a calibration/reproducibility stage (EXP1) followed by a primary comparative/robustness benchmark (EXP2); the current merged recovery snapshot contains 299 committed result artifacts (99.7\% artifact coverage) plus a 30-row SHAP recovery batch, yielding 275 analyzable unique runs out of 300 planned cells (91.7\%). Across complete model-size blocks ($5$ models, $N \in \{50,100,200\}$), Friedman tests indicate significant method differences for fidelity ($\chi^2=42.12$, $p=3.78\times10^{-9}$), stability ($\chi^2=40.68$, $p=7.65\times10^{-9}$), sparsity ($\chi^2=35.64$, $p=8.92\times10^{-8}$), faithfulness gap ($\chi^2=45.00$, $p=9.25\times10^{-10}$), and runtime ($\chi^2=30.44$, $p=1.12\times10^{-6}$). SHAP leads on fidelity/stability, DiCE leads on sparsity, and LIME remains fastest overall. We release the framework, operation protocol, and artifacts with explicit data-quality caveats for reproducible benchmark use under a quantitative-only claim scope.
